\chapter{Criterio de mapa público mínimo del corpus}

\section{Mapa público mínimo}
\begin{definition}
Se llama mapa público mínimo del corpus a la cartografía externa más sobria mediante la cual pueden presentarse, de forma legible y disciplinada, las principales piezas publicables del sistema sin abrir todavía el corpus completo ni borrar su jerarquía metodológica.
\end{definition}

\begin{observation}
El mapa público mínimo no sustituye ni a la tabla madre ni al dossier externo. Su función es ofrecer una visión externa condensada del conjunto, manteniendo visibles estructura, niveles y tipos de piezas sin inflar el alcance público del sistema.
\end{observation}

\section{Condición mínima de mapeabilidad pública}
\begin{definition}
Existe mapa público mínimo del corpus cuando pueden declararse simultáneamente:
\begin{enumerate}
  \item una selección suficientemente sobria de piezas publicables;
  \item una organización legible entre claims, fichas, dossiers y reservas;
  \item vínculo visible con la genealogía metodológica interna;
  \item distinción suficiente entre núcleo, material observado y capa pública;
  \item ausencia de compresión abusiva del sistema completo en una sola figura externa.
\end{enumerate}
\end{definition}

\begin{observation}
No se exige todavía una exposición integral de todos los materiales. Basta con una cartografía mínima que permita orientarse sin falsear la arquitectura real del corpus.
\end{observation}

\section{Mapa y ficha pública}
\begin{definition}
Toda ficha pública mínima debe poder ubicarse dentro del mapa público mínimo como unidad externa identificable y no como fragmento aislado sin contexto.
\end{definition}

\begin{observation}
Esta cláusula permite que la pieza individual conserve relación con el conjunto sin necesidad de exponer todo el sistema en cada ficha.
\end{observation}

\section{Mapa y dossier}
\begin{definition}
Todo dossier externo de claims seleccionados debe poder ubicarse dentro del mapa público mínimo como agrupación legible, con función reconocible y límites sobrios respecto del resto del corpus.
\end{definition}

\begin{observation}
El mapa público no debe disolver la identidad del dossier ni convertirlo en una simple lista plana de claims.
\end{observation}

\section{Condición mínima de no mapeabilidad}
\begin{definition}
No debe constituirse mapa público mínimo cuando la cartografía externa dependa de simplificaciones que destruyan la diferencia entre claim, dossier, reserva, fundamento y material todavía interno.
\end{definition}

\begin{observation}
La no mapeabilidad pública no implica fracaso editorial. Puede significar solo que el sistema necesita una mediación más sobria antes de exponerse como conjunto.
\end{observation}

\section{Cláusula de alcance}
\begin{observation}
El criterio de mapa público mínimo del corpus fija una primera cartografía externa sobria del sistema. No constituye todavía una presentación pública integral, definitiva o exhaustiva del conjunto MADELEINE.
\end{observation}

